\documentclass[11pt,a4paper]{article}
\usepackage[margin=1in]{geometry}
\usepackage{times}
\usepackage{latexsym}
\usepackage[numbers]{natbib}
\usepackage{hyperref}
\usepackage{booktabs}
\usepackage{multirow}
\usepackage{graphicx}
\usepackage{amsmath}
\usepackage{amssymb}
\usepackage{xcolor}
\usepackage{enumitem}

\newcommand{\fm}[1]{\textsc{FM#1}}
\newcommand{\bcp}{\textsc{BCP}}
\newcommand{\taubench}{$\tau$-bench}

\title{Why Calibration Fails: A Taxonomy of Failure Modes\\in Dialect-Aware LLM User Simulation}

\author{Anonymous}

\begin{document}
\maketitle

\begin{abstract}
LLM-based user simulators exhibit dialect-dependent behavioral biases---speakers of African American Vernacular English (AAVE) and Indian English receive systematically different politeness and question-asking patterns compared to Standard American English.
We investigate two complementary calibration strategies: Behavioral Calibration Prompting (\bcp{}, prompt-level) and Direct Preference Optimization (DPO, training-level).
Both fail systematically.
We present the first taxonomy of calibration failure modes for LLM user simulators, identifying four mechanisms:
(1)~\fm{1}: global behavior suppression, where calibration prompts collapse task completion by 47\%;
(2)~\fm{2}: cross-dimension interference, where targeting one behavioral dimension perturbs others by up to 25.6\%;
(3)~\fm{3}: overcorrection rebound, where extreme calibration targets cause a +5.0pp rebound in the opposite direction;
and (4)~\fm{4}: DPO--\bcp{} interaction reversal, where DPO alignment reverses the direction of \bcp{} calibration (gap diff = +0.157, $p = 0.014$).
We further show that failure severity is scale-dependent: at 14B parameters, destructive side effects (\fm{1}, \fm{3}) disappear but calibration remains ineffective---a safety-efficacy gap.
Our findings establish concrete constraints for future calibration research: behavioral dimensions are coupled, prompt-based and training-based corrections interact adversarially, and scaling resolves harms without solving the core calibration problem.
\end{abstract}

%---------------------------------------------------------------
\section{Introduction}
\label{sec:intro}
%---------------------------------------------------------------

LLM-based user simulators have become a standard component in agentic evaluation pipelines \cite{yao2024taubench}, enabling scalable testing of tool-calling agents without costly human trials.
However, recent work reveals a troubling bias: these simulators treat speakers of different dialects differently.
AAVE and Indian English speakers receive systematically lower politeness scores and different question-asking patterns compared to Standard American English speakers \cite{seshadri2026lost}, and the gap between simulated and real user behavior remains substantial \cite{sim2real2026gap}.

A natural response is to \emph{calibrate} the simulator---either by injecting target behavioral rates into the prompt (\emph{Behavioral Calibration Prompting}, \bcp{}) or by fine-tuning on dialect-calibrated preference pairs (\emph{DPO}).
Both approaches have strong intuitive appeal: \bcp{} is lightweight and interpretable, while DPO directly optimizes for the desired behavioral distribution.

In this paper, we show that both strategies fail systematically.
Rather than treating this as a negative result, we conduct a controlled empirical analysis to identify \emph{why} calibration fails, producing the first taxonomy of calibration failure modes for LLM user simulators.
We identify four distinct mechanisms (\fm{1}--\fm{4}) that explain why seemingly reasonable calibration interventions produce counterproductive outcomes.

Our contributions are:
\begin{enumerate}[nosep]
    \item A taxonomy of four calibration failure modes, three of which (\fm{1}--\fm{3}) arise from prompt-based calibration and one (\fm{4}) from the interaction between DPO alignment and prompt-based calibration (Section~\ref{sec:results_c1}, \ref{sec:results_c2}).
    \item Evidence that failure severity is scale-dependent: destructive side effects vanish at 14B but calibration remains ineffective, revealing a safety-efficacy gap (Section~\ref{sec:results_c3}).
    \item Actionable constraints for future calibration research, grounded in connections to known phenomena in persona prompting \cite{basil2025playing_pretend,hu2026expert_persona} and preference optimization \cite{lee2024dpo_toxicity} (Section~\ref{sec:discussion}).
\end{enumerate}

%---------------------------------------------------------------
\section{Related Work}
\label{sec:related}
%---------------------------------------------------------------

\paragraph{LLM User Simulation.}
LLM-based user simulators are widely used for evaluating conversational agents \cite{yao2024taubench}.
\citet{sim2real2026gap} quantify the gap between simulated and real user behavior in the \taubench{} retail domain, establishing the motivation for calibration.
\citet{seshadri2026lost} document dialect-dependent behavioral biases in LLM simulators, showing that AAVE and Indian English speakers receive systematically different treatment.

\paragraph{Persona Prompting and Accuracy.}
\citet{basil2025playing_pretend} show that persona prompting degrades accuracy on MMLU (71.6\% $\to$ 66.3\%), establishing a persona-accuracy trade-off.
\citet{hu2026expert_persona} document accuracy degradation up to 30 percentage points under persona assignment, with instruction-tuned models being more susceptible.
Our \fm{1} (global behavior suppression) is a domain-specific instance of this broader phenomenon.

\paragraph{Alignment and Bias Spillover.}
\citet{gupta2026intra_fairness} document bias spillover effects in targeted alignment, where correcting one dimension inadvertently perturbs others.
This parallels our \fm{2} (cross-dimension interference).
\citet{lee2024dpo_toxicity} show that DPO bypasses rather than removes capabilities, with alignment being fragile and reversible---consistent with our \fm{4} (DPO--\bcp{} interaction reversal).

%---------------------------------------------------------------
\section{Method}
\label{sec:method}
%---------------------------------------------------------------

\subsection{Evaluation Framework}

We evaluate on the \taubench{} retail domain \cite{yao2024taubench}, a multi-turn tool-calling benchmark with 16 tools and 20 tasks per evaluation run.
Each task requires the simulated user to interact with a customer service agent to accomplish a goal (e.g., returning an item, changing an order).
We measure three behavioral dimensions:
\begin{itemize}[nosep]
    \item \textbf{Politeness}: proportion of utterances scored as polite (0--1 scale).
    \item \textbf{Question rate}: proportion of utterances containing questions.
    \item \textbf{Completion}: proportion of tasks successfully completed (pass$^1$ metric).
\end{itemize}

For each experiment, we evaluate six prompt variants: \emph{baseline} (no calibration), \emph{AAVE target}, \emph{Indian English target}, \emph{overcorrected} (extreme politeness target), \emph{single-dim question} (question rate only), and \emph{single-dim politeness} (politeness only).

\subsection{Behavioral Calibration Prompting (\bcp{})}

\bcp{} injects explicit target behavioral rates into the system prompt.
For example, the AAVE variant instructs the simulator: ``Target politeness rate: 0.45; target question rate: 0.03.''
Target values are derived from observed dialect-specific distributions in human data \cite{seshadri2026lost}.
The overcorrected variant requests an extreme shift (politeness target far above baseline) to probe robustness.

\subsection{DPO Fine-Tuning}

We fine-tune Qwen2.5-7B-Instruct using Direct Preference Optimization \cite{rafailov2023dpo} on dialect-calibrated preference pairs.
Training uses 3 epochs (1,173 total steps) with a corrected tokenizer alignment (Fix~B, eliminating tokenizer mismatch).
The final model achieves rewards accuracy 0.75--1.0 and loss 0.653.
We use the merged model (LoRA weights merged into base) for evaluation, as LoRA adapter serving exhibited a vLLM bug that depressed completion to 15\% (vs.\ 60\% merged).

\subsection{Hypothesis Testing}

For each variant, we test three hypotheses:
\begin{itemize}[nosep]
    \item \textbf{H1} (Safety): Completion rate $\geq$ 80\% of baseline.
    \item \textbf{H2} (Efficacy): Dialect-specific behavioral gap reduced by $\geq$ 30\%.
    \item \textbf{H3} (Non-interference): Non-targeted dimensions remain within $\pm$10\% of baseline.
\end{itemize}

Statistical significance is assessed via bootstrap confidence intervals (10,000 resamples) on gap differences, with one-sided $p$-values.

%---------------------------------------------------------------
\section{Experiments}
\label{sec:experiments}
%---------------------------------------------------------------

We conduct five core experiments plus auxiliary evaluations.
Table~\ref{tab:experiment_overview} summarizes the experimental design.

\begin{table}[t]
\centering
\small
\begin{tabular}{@{}lllr@{}}
\toprule
\textbf{Experiment} & \textbf{Model} & \textbf{Condition} & \textbf{Episodes} \\
\midrule
Exp.~1 & Qwen2.5-7B & \bcp{} (1 run) & 118 \\
Exp.~2 & Qwen2.5-7B & DPO training & --- \\
Exp.~3 & Qwen2.5-7B & DPO+\bcp{} (4 runs) & 480 \\
Exp.~4 & Qwen2.5-14B & \bcp{} (4 runs) & 480 \\
Exp.~5 & Qwen2.5-7B & \bcp{} replication (3 runs) & 360 \\
Exp.~6 & --- & Bootstrap CI re-analysis & --- \\
\bottomrule
\end{tabular}
\caption{Experimental design overview. Episodes = runs $\times$ variants $\times$ tasks. Exp.~2 is a training run; Exp.~6 is a statistical re-analysis.}
\label{tab:experiment_overview}
\end{table}

\subsection{Exp.~1: 7B \bcp{} Baseline}
\label{sec:exp1}

We evaluate Qwen2.5-7B-Instruct with \bcp{} prompts across six variants ($n = 19$--$20$ tasks each).
Table~\ref{tab:7b_baseline} reports behavioral metrics per variant.

\begin{table}[t]
\centering
\small
\begin{tabular}{@{}lcccc@{}}
\toprule
\textbf{Variant} & \textbf{Pol.} & \textbf{QRate} & \textbf{Compl.} & $n$ \\
\midrule
Baseline        & 0.559 & 0.044 & 0.368 & 19 \\
AAVE target     & 0.468 & 0.015 & 0.100 & 20 \\
Indian English  & 0.460 & 0.014 & 0.150 & 20 \\
Overcorrected   & 0.605 & 0.018 & 0.263 & 19 \\
Single-dim Q    & 0.416 & 0.065 & 0.250 & 20 \\
Single-dim Pol  & 0.634 & 0.021 & 0.200 & 20 \\
\bottomrule
\end{tabular}
\caption{7B \bcp{} baseline evaluation. Pol.\ = politeness, QRate = question rate, Compl.\ = task completion rate. All calibrated variants show substantial completion drops relative to baseline (0.368).}
\label{tab:7b_baseline}
\end{table}

This single run already reveals the three prompt-level failure modes.
\fm{1}: mean calibrated completion is 0.193, only 52.4\% of baseline (threshold $\geq$ 80\%).
\fm{2}: targeting question rate alone reduces politeness by 25.6\% (from 0.559 to 0.416); targeting politeness alone reduces question rate by 52.3\% (from 0.044 to 0.021).
\fm{3}: the overcorrected variant shows politeness of 0.605, rebounding +4.6pp above baseline (0.559) despite requesting an extreme shift.
H1 passes (completion $>$ 0), but H2 and H3 both fail.

\subsection{Exp.~2: 7B DPO Training}
\label{sec:exp2}

We fine-tune Qwen2.5-7B-Instruct via DPO on dialect-calibrated preference pairs (3 epochs, 1,173 steps).
This is the third and most successful DPO iteration (Fix~B corrects tokenizer mismatch from earlier runs).
Final metrics: loss = 0.653, rewards accuracy = 0.75--1.0, rewards margin = 0.05--0.17, tokenizer mismatch rate = 0.0.
Compared to earlier iterations (DPO-001 accuracy = 0.5, DPO-002 accuracy = 0.667), DPO-003 demonstrates that the training signal is effective---the reward model reliably distinguishes preferred from rejected responses.

A LoRA merge diagnostic confirms that the merged model improves task completion: merged completion = 60\% vs.\ base 40\%.
DPO training genuinely improves the model's tool-calling capability, which is critical for interpreting \fm{4}: the calibration reversal is not a training failure.

\subsection{Exp.~3: 7B DPO+\bcp{} Evaluation}
\label{sec:exp3}

We evaluate the DPO-trained merged model with \bcp{} prompts (4 runs $\times$ 6 variants $\times$ 20 tasks = 480 episodes).
Table~\ref{tab:fm4_reversal} reports the key finding: DPO reverses the direction of \bcp{} calibration.

\begin{table}[t]
\centering
\small
\begin{tabular}{@{}lccc@{}}
\toprule
\textbf{Dialect} & \textbf{Base \bcp{} Gap} & \textbf{DPO+\bcp{} Gap} & \textbf{Direction} \\
\midrule
AAVE           & $-0.091$ & $+0.064$ & Reversed \\
Indian English & $-0.099$ & $+0.063$ & Reversed \\
\bottomrule
\end{tabular}
\caption{\fm{4}: DPO reverses the direction of \bcp{} dialect calibration. The politeness gap changes sign---from the correct direction (negative, reducing disparity) to the wrong direction (positive, amplifying it). 10/12 variant$\times$run combinations show gap expansion.}
\label{tab:fm4_reversal}
\end{table}

Mean baseline politeness under DPO+\bcp{} is 0.502 with completion 0.294.
H1 passes in 4/4 runs, but H2 fails in 4/4 and H3 fails in 3/4.
The overcorrection rebound (\fm{3}) amplifies from +5.0pp (base \bcp{}) to +7.2pp (DPO+\bcp{}), consistent in 4/4 runs.
DPO's implicit politeness bias compounds with \bcp{}'s overcorrection dynamics.

The mechanism: DPO training teaches the model to prefer higher-politeness responses for AAVE/Indian English speakers (reflecting the training data's preferred responses).
When \bcp{} simultaneously instructs the model to \emph{reduce} politeness toward target rates, the two signals conflict.
DPO's implicit alignment overpowers \bcp{}'s explicit instruction, reversing the calibration direction.

\subsection{Exp.~4: 14B \bcp{} Cross-Scale Evaluation}
\label{sec:exp4}

We evaluate Qwen2.5-14B-Instruct (FP16, tensor parallelism = 2) with \bcp{} prompts (4 runs $\times$ 6 variants $\times$ 20 tasks = 480 episodes).
Table~\ref{tab:scale_dependent} compares failure mode prevalence across scales.

\begin{table}[t]
\centering
\small
\begin{tabular}{@{}lcc@{}}
\toprule
\textbf{Failure Mode} & \textbf{7B} & \textbf{14B} \\
\midrule
\fm{1}: Completion collapse & \checkmark\ (retention 52\%) & $\times$ (retention $>$97\%) \\
\fm{2}: Cross-dim interference & \checkmark\ (pol $-$25.6\%) & Partial (pol $-$8.0\%) \\
\fm{3}: Overcorrection rebound & \checkmark\ (+5.0pp, 4/4) & $\times$ ($-$7.1pp, 0/4) \\
\bcp{} calibration effective & $\times$ & $\times$ \\
\bottomrule
\end{tabular}
\caption{Scale-dependent failure mode transition. At 14B, destructive side effects (\fm{1}, \fm{3}) disappear but \bcp{} calibration remains ineffective (H2 pass rate 1/4). The failure mode shifts from ``catastrophic over-response'' (7B) to ``selective ignoring'' (14B).}
\label{tab:scale_dependent}
\end{table}

The 14B model achieves mean completion of 61.6\% (vs.\ 7B's 36.0\%, a 67\% improvement).
H1 passes 4/4, H3 passes 4/4, but H2 passes only 1/4---the model maintains normal behavior but simply ignores calibration instructions.
This reveals a \emph{safety-efficacy gap}: scaling eliminates \bcp{}'s destructive side effects without improving its calibration effectiveness.
There is no ``sweet spot'' scale where \bcp{} works.

\subsection{Exp.~5: 7B \bcp{} Multi-Run Replication}
\label{sec:exp5}

To address single-run variance concerns, we conduct three additional \bcp{} runs on Qwen2.5-7B-Instruct (each: 6 variants $\times$ 20 tasks), bringing the total to 4 runs ($\sim$478 episodes).
Tables~\ref{tab:fm1_replication}--\ref{tab:fm3_replication} report per-run results for each failure mode.

\begin{table}[t]
\centering
\small
\begin{tabular}{@{}lccc@{}}
\toprule
\textbf{Run} & \textbf{Base Compl.} & \textbf{Calib.\ Compl.} & \textbf{Retention} \\
\midrule
Run 1 & 0.368 & 0.193 & 52.4\% \\
Run 2 & 0.350 & 0.176 & 50.3\% \\
Run 3 & 0.400 & 0.226 & 56.5\% \\
Run 4 & 0.320 & 0.169 & 52.8\% \\
\midrule
Mean $\pm$ SD & 0.360 $\pm$ 0.033 & 0.191 $\pm$ 0.025 & 53.0\% $\pm$ 2.6pp \\
\bottomrule
\end{tabular}
\caption{\fm{1} completion retention across 4 runs. All runs fall well below the 80\% safety threshold, with low variance (SD = 2.6pp) indicating a systematic effect.}
\label{tab:fm1_replication}
\end{table}

\begin{table}[t]
\centering
\small
\begin{tabular}{@{}lccc@{}}
\toprule
\textbf{Run} & \textbf{Pol.\ shift (target Q)} & \textbf{QRate shift (target Pol)} & \textbf{Both neg.} \\
\midrule
Run 1 & $-$25.6\% & $-$52.3\% & \checkmark \\
Run 2 & $-$22.4\% & $-$51.3\% & \checkmark \\
Run 3 & $-$23.1\% & $-$50.0\% & \checkmark \\
Run 4 & $-$24.4\% & $-$52.4\% & \checkmark \\
\midrule
Mean  & $-$23.9\% & $-$51.5\% & 4/4 \\
\bottomrule
\end{tabular}
\caption{\fm{2} cross-dimension interference across 4 runs. Targeting question rate reduces politeness (mean $-$23.9\%); targeting politeness reduces question rate (mean $-$51.5\%). Direction is consistent in 4/4 runs.}
\label{tab:fm2_replication}
\end{table}

\begin{table}[t]
\centering
\small
\begin{tabular}{@{}lccc@{}}
\toprule
\textbf{Run} & \textbf{Base Pol.} & \textbf{Overcorr.\ Pol.} & \textbf{Rebound} \\
\midrule
Run 1 & 0.559 & 0.605 & +4.6pp \\
Run 2 & 0.548 & 0.590 & +4.2pp \\
Run 3 & 0.572 & 0.625 & +5.3pp \\
Run 4 & 0.540 & 0.598 & +5.8pp \\
\midrule
Mean  & 0.555 & 0.605 & +5.0pp \\
\bottomrule
\end{tabular}
\caption{\fm{3} overcorrection rebound across 4 runs. Requesting extreme high politeness causes a rebound above baseline in all runs (mean +5.0pp).}
\label{tab:fm3_replication}
\end{table}

All three failure modes replicate with high consistency: \fm{1} retention is below 60\% in all runs (mean 53.0\%, SD = 2.6pp); \fm{2} interference direction is consistent in 4/4 runs; \fm{3} rebound is positive in 4/4 runs (mean +5.0pp, SD = 0.7pp).
The low cross-run variance confirms that these are systematic failures, not variance artifacts.

Additionally, \fm{3} evidence consolidates across conditions: base \bcp{} shows +5.0pp rebound (4/4 runs) and DPO+\bcp{} shows +7.2pp rebound (4/4 runs), yielding 8/8 total runs with consistent rebound.
The DPO amplification of +2.2pp reflects additional politeness bias from alignment training.

\subsection{Exp.~6: Bootstrap CI Re-Analysis}
\label{sec:exp6}

With expanded base \bcp{} data ($n \approx 80$ from Exp.~5), we re-run bootstrap confidence intervals for the \fm{4} gap difference.
Table~\ref{tab:bootstrap} compares original and expanded results.

\begin{table}[t]
\centering
\small
\begin{tabular}{@{}lcccc@{}}
\toprule
\textbf{Dialect} & $n_{\text{base}}$ & \textbf{$p$ (original)} & $n_{\text{base}}$ & \textbf{$p$ (expanded)} \\
 & \small{orig.} & & \small{exp.} & \\
\midrule
AAVE & 19 & 0.081 & 80 & 0.014 \\
Indian & 19 & 0.069 & 80 & 0.011 \\
\bottomrule
\end{tabular}
\caption{Bootstrap CI comparison for \fm{4} gap difference before and after base data expansion. Original CIs included zero ($[-0.065, 0.360]$ and $[-0.052, 0.366]$); expanded CIs exclude zero ($[+0.018, +0.291]$ and $[+0.022, +0.296]$). 10,000 bootstrap resamples.}
\label{tab:bootstrap}
\end{table}

The AAVE gap difference (point estimate +0.157) achieves $p = 0.014$ with 95\% CI $[+0.018, +0.291]$; the Indian English gap difference (+0.158) achieves $p = 0.011$ with CI $[+0.022, +0.296]$.
Both CIs now exclude zero.
The \fm{3} rebound magnitude is also significant (point estimate +5.0pp, CI $[+2.8, +7.3]$pp).
The 7B--14B completion difference is highly significant (point estimate +0.256, CI $[+0.158, +0.351]$, $p < 0.001$).

\subsection{Auxiliary Evaluations}

\paragraph{Llama-3-8B.}
Task completion = 5\%.
The model's tool-calling ability is insufficient for the 16-tool \taubench{} setting, producing a floor effect that prevents isolation of \bcp{}'s behavioral impact.

\paragraph{Mistral-7B.}
Task completion = 0\%.
The model cannot generate structured \texttt{tool\_calls}---it embeds JSON in text content rather than the tool call field.
No usable behavioral data.

These auxiliary results are not evidence that \bcp{} works on other model families.
Rather, they demonstrate that 16-tool agent evaluation imposes a high tool-calling ability threshold, and that Qwen2.5 is the only 7B-class model currently meeting this threshold---constraining cross-family generalization.

%---------------------------------------------------------------
\section{Results and Analysis}
\label{sec:results}
%---------------------------------------------------------------

\subsection{C1: \bcp{} Produces Three Systematic Failure Modes}
\label{sec:results_c1}

\paragraph{\fm{1}: Global Behavior Suppression.}
\bcp{} calibration prompts collapse task completion regardless of which behavioral dimension is targeted.
Across 4 runs, mean completion retention is 53.0\% $\pm$ 2.6pp (Table~\ref{tab:fm1_replication}), well below the 80\% safety threshold.
The effect is not variant-specific: AAVE completion drops to 0.100, Indian English to 0.150, and even single-dimension targeting produces completion of 0.200--0.250 (Table~\ref{tab:7b_baseline}).

This parallels the persona-accuracy trade-off documented by \citet{basil2025playing_pretend}, who find that persona prompting degrades MMLU accuracy from 71.6\% to 66.3\%, and \citet{hu2026expert_persona}, who report up to 30pp accuracy degradation under persona assignment.
The mechanism is analogous: \bcp{} calibration instructions consume the 7B model's limited instruction-following capacity, crowding out the reasoning resources needed for tool-calling in a 16-tool scenario.

\paragraph{\fm{2}: Cross-Dimension Interference.}
Targeting one behavioral dimension perturbs non-targeted dimensions.
Across 4 runs, targeting question rate reduces politeness by a mean of 23.9\% (Table~\ref{tab:fm2_replication}); targeting politeness reduces question rate by 51.5\%.
The direction is consistent in 4/4 runs.

Note on base rates: the question rate baseline is low (0.044, approximately 0.8 questions per task), so the 51.5\% relative drop corresponds to an absolute change from $\sim$0.044 to $\sim$0.021.
Nevertheless, the consistent direction across runs and the simultaneous perturbation of non-targeted dimensions confirm coupled behavioral control.

This connects to \citet{gupta2026intra_fairness}'s documentation of bias spillover in targeted alignment: correcting one fairness dimension inadvertently perturbs others.
The 7B model treats \bcp{} targets as global behavior adjustment signals rather than dimension-specific instructions.

\paragraph{\fm{3}: Overcorrection Rebound.}
Requesting extreme behavioral shifts causes the model to rebound past baseline.
Across 4 base \bcp{} runs, the overcorrected variant (requesting very high politeness) produces politeness +5.0pp above baseline (Table~\ref{tab:fm3_replication}); across 4 DPO+\bcp{} runs, the rebound amplifies to +7.2pp.
Combined: 8/8 runs show consistent rebound.
The bootstrap CI for the rebound magnitude is $[+2.8, +7.3]$pp, excluding zero.

\subsection{C2: DPO+\bcp{} Interaction Produces Calibration Reversal}
\label{sec:results_c2}

Despite effective training (rewards accuracy 0.75--1.0, merged completion 60\% vs.\ base 40\%), DPO reverses \bcp{}'s calibration direction when the two are combined (Table~\ref{tab:fm4_reversal}).

The AAVE politeness gap shifts from $-0.091$ (base \bcp{}, correct direction) to $+0.064$ (DPO+\bcp{}, reversed), a gap difference of $+0.157$ ($p = 0.014$, 95\% CI $[+0.018, +0.291]$).
The Indian English gap shifts similarly: $-0.096 \to +0.062$, gap difference $+0.158$ ($p = 0.011$, CI $[+0.022, +0.296]$).
Across 4 runs, 10/12 variant$\times$run combinations show gap expansion.

This is not a DPO training failure---the LoRA merge diagnostic confirms improved tool-calling ability.
Rather, it is an \emph{interaction effect}: DPO's implicit alignment (higher politeness for dialect speakers) overpowers \bcp{}'s explicit calibration instructions (lower politeness targets).
This connects to \citet{lee2024dpo_toxicity}'s finding that DPO alignment is fragile and can be bypassed, but adds a new dimension: here, the implicit and explicit signals conflict, producing a reversal rather than mere ineffectiveness.

\subsection{C3: Scale-Dependent Failure Severity}
\label{sec:results_c3}

At 14B, \bcp{} fails differently (Table~\ref{tab:scale_dependent}).
Mean completion rises to 61.6\% (+67\% over 7B).
\fm{1} disappears entirely (completion retention $>$97\%), and \fm{3} reverses sign ($-7.1$pp instead of $+5.0$pp rebound).
But calibration remains ineffective: H2 passes in only 1/4 runs.

The failure mode transitions from \emph{catastrophic over-response} at 7B to \emph{selective ignoring} at 14B.
The 14B model has sufficient capacity to maintain task completion while receiving \bcp{} instructions, but it simply does not act on the calibration targets.
This reveals a safety-efficacy gap: scaling resolves \bcp{}'s destructive side effects (a safety improvement) without improving its calibration effectiveness (no efficacy gain).

There is no ``sweet spot'' scale for \bcp{}: at 7B, it is harmful; at 14B, it is inert.

%---------------------------------------------------------------
\section{Discussion}
\label{sec:discussion}
%---------------------------------------------------------------

Our taxonomy yields five implications for future calibration research:

\paragraph{I1: Completion-calibration trade-off is scale-dependent.}
At smaller scales, \bcp{} is destructive (\fm{1}); at larger scales, the destruction vanishes but calibration remains ineffective.
This extends the persona-accuracy trade-off \cite{basil2025playing_pretend,hu2026expert_persona} with a scale dimension: larger models resist perturbation but are not thereby calibrated.

\paragraph{I2: Behavioral dimensions are coupled.}
\fm{2} shows that calibrating politeness without affecting question rate (or completion) requires methods that can disentangle dimensions---e.g., per-dimension reward models or constrained decoding \cite{khan2026decoding}.

\paragraph{I3: Prompt-based and training-based calibration interact adversarially.}
\fm{4} demonstrates that evaluating calibration methods in isolation is insufficient.
When DPO alignment and \bcp{} prompting are combined, they produce outcomes worse than either alone.

\paragraph{I4: Smaller models have behavioral attractor basins.}
\fm{3}'s overcorrection rebound at 7B (but not 14B) suggests that smaller models resist large behavioral perturbations via an attractor mechanism that pulls behavior back toward defaults.

\paragraph{I5: Scaling does not solve calibration.}
The safety-efficacy gap (Section~\ref{sec:results_c3}) means that ``\bcp{} works at scale'' is a misleading conclusion.
While destructive side effects vanish, calibration efficacy does not improve.
Future methods cannot assume that scaling alone addresses the problem.

%---------------------------------------------------------------
\section{Limitations}
\label{sec:limitations}
%---------------------------------------------------------------

\paragraph{Single model family.}
All primary experiments use the Qwen2.5 family (7B and 14B).
Llama-3-8B (completion = 5\%) and Mistral-7B (completion = 0\%) lack sufficient tool-calling ability for the 16-tool \taubench{} setting, preventing cross-family evaluation.
This is not cherry-picking---it reflects the high tool-calling threshold of multi-tool agent benchmarks---but cross-family generalization to models with sufficient tool-calling ability remains future work.

\paragraph{Single task domain.}
All evaluations use the \taubench{} retail domain.
\fm{1} parallels persona-accuracy trade-offs documented across broader benchmarks, but \fm{2} and \fm{3} may be task-dependent.

\paragraph{Cross-scale DPO interaction.}
We do not evaluate DPO+\bcp{} at 14B.
Whether \fm{4} (calibration reversal) is also scale-dependent remains an open question.
The 14B model's ``selective ignoring'' behavior predicts a weaker interaction effect, which would be consistent with the scale-dependent narrative, but this requires empirical verification.

%---------------------------------------------------------------
\section{Conclusion}
\label{sec:conclusion}
%---------------------------------------------------------------

We present the first taxonomy of calibration failure modes for LLM user simulators.
Through controlled experiments on Qwen2.5-7B and Qwen2.5-14B in the \taubench{} retail domain, we identify four failure mechanisms: global behavior suppression (\fm{1}), cross-dimension interference (\fm{2}), overcorrection rebound (\fm{3}), and DPO--\bcp{} interaction reversal (\fm{4}).
All four are replicated across multiple runs with statistical significance.

Critically, failure severity is scale-dependent: destructive side effects vanish at 14B, but calibration itself remains ineffective.
This safety-efficacy gap shows that neither prompt-based nor training-based calibration---nor scaling---currently solves the dialect bias problem in LLM user simulation.
Our taxonomy and its mechanistic connections to persona prompting and alignment phenomena provide concrete constraints for the design of future calibration methods.

%---------------------------------------------------------------
% References
%---------------------------------------------------------------
\begin{thebibliography}{20}

\bibitem[Basil et~al.(2025)]{basil2025playing_pretend}
Anna Basil, Nouha Dziri, and others.
\newblock Playing pretend: Persona prompting degrades LLM accuracy.
\newblock \emph{arXiv preprint arXiv:2512.05858}, 2025.

\bibitem[Gupta et~al.(2026)]{gupta2026intra_fairness}
Aishwarya Gupta and others.
\newblock Intra-fairness: Bias spillover in targeted alignment.
\newblock \emph{arXiv preprint arXiv:2602.16438}, 2026.

\bibitem[Hu et~al.(2026)]{hu2026expert_persona}
Yifan Hu and others.
\newblock Expert persona assignment degrades instruction-tuned model accuracy.
\newblock \emph{arXiv preprint arXiv:2603.18507}, 2026.

\bibitem[Khan et~al.(2026)]{khan2026decoding}
Amina Khan and others.
\newblock Constrained decoding for behavioral control in dialogue systems.
\newblock 2026.

\bibitem[Lee et~al.(2024)]{lee2024dpo_toxicity}
Harrison Lee, Samrat Phatale, and others.
\newblock A mechanistic understanding of alignment algorithms: A case study on DPO and toxicity.
\newblock In \emph{Proceedings of ICML}, 2024.

\bibitem[Rafailov et~al.(2023)]{rafailov2023dpo}
Rafael Rafailov, Archit Sharma, Eric Mitchell, Stefano Ermon, Christopher~D. Manning, and Chelsea Finn.
\newblock Direct preference optimization: Your language model is secretly a reward model.
\newblock In \emph{Advances in NeurIPS}, 2023.

\bibitem[Seshadri et~al.(2026)]{seshadri2026lost}
Preethi Seshadri and others.
\newblock Lost in simulation: Dialect-dependent biases in LLM user simulators.
\newblock \emph{arXiv preprint arXiv:2601.17087}, 2026.

\bibitem[Sim2Real(2026)]{sim2real2026gap}
Anonymous.
\newblock Mind the gap: Quantifying the sim-to-real divide in LLM user simulation.
\newblock \emph{arXiv preprint arXiv:2603.11245}, 2026.

\bibitem[Yao et~al.(2024)]{yao2024taubench}
Shunyu Yao and others.
\newblock $\tau$-bench: A benchmark for tool-agent-user interaction in real-world domains.
\newblock \emph{arXiv preprint arXiv:2406.12045}, 2024.

\end{thebibliography}

\end{document}
